\section{Conclusion}
\label{sec:conclusion}

Section~\ref{sec:introduction} opened with a gap: quantum-inspired
evolutionary algorithms have been applied to multi-objective problems
since at least \citet{kim2006qmea}'s quantum-inspired multi-objective
knapsack algorithm, yet no prior survey had taxonomized or critically
compared QIEA-for-MOO methods on their own terms, separately from
single-objective QIEA work on one side and classical MOO on the other.
This survey addressed that gap directly. Section~\ref{sec:taxonomy}
organized the literature identified in
Section~\ref{sec:introduction:scope} along five axes -- representation
scheme, quantum-inspired operator type, selection and diversity
mechanism, hybridization with classical MOEAs, and application domain.
Section~\ref{sec:feature-comparison} applied those axes systematically
across the surveyed methods, and in doing so surfaced a pattern the
taxonomy alone had only implied: an overwhelming concentration on
NSGA-II's crowding distance as the selection mechanism of choice, and a
benchmarking culture in which every method is compared against that
same classical baseline but never against its own QIEA-MOO siblings.
Section~\ref{sec:critical-analysis} traced that pattern, and several
others like it, into per-category weaknesses and three cross-cutting
problems -- reproducibility gaps, weak benchmarking practices, and
scalability claims rarely stress-tested -- that recur across the field
regardless of which taxonomy category a given method falls into.

Section~\ref{sec:case-study} grounded several of these abstract claims
in a concrete worked example: this project's own block-partitioned
quantum-circuit-optimization pipeline, itself a direct extension of
\citet{ghlib2026scalable}'s scalable, block-partitioned NSGA-II
framework. We want to restate plainly what that case study was and was
not. It was illustrative -- a way to show, with a real,
actively-maintained codebase, what a swappable selection mechanism
looks like in practice, how an adaptive hypervolume computation can
silently mislead a cross-algorithm comparison, and how easily a
reproducibility gap can hide even inside careful, tested development.
It was not this survey's primary evidence, and none of its numbers
should be read as a benchmark result for QIEA-MOO methods themselves,
which the case study's own pipeline does not use. Full empirical
validation of that pipeline -- including the mutation-scheme and
hybrid-local-search ablations Section~\ref{sec:case-study:mutation}
explicitly left unresolved, pending re-verification on the pipeline's
corrected code -- is the responsibility of a separate, forthcoming
empirical results paper, not this one.

Section~\ref{sec:open-challenges} closed by turning the critical
analysis forward rather than only backward: untested pairings between
the Q-bit representation and selection mechanisms beyond NSGA-II's,
a still-undemonstrated interference-based operator, concrete extensions
to the case-study pipeline drawn from adjacent quantum-computing
literature, and the wider field's movement toward learned and
neural-guided search that neither QIEA-MOO nor this survey's case study
has yet engaged with. None of these close the gap this survey opened
with; each narrows it. If the motivation from
Section~\ref{sec:introduction} is taken seriously -- that a cheaper,
shallower optimization process is not merely a performance metric but
part of a broader case for energy-efficient
computing~\citep{banciu2026aiot,moga2025rlmicrogrid} -- then QIEA-MOO's
comparative youth as a subfield, documented throughout this survey, is
better read as room to grow into than as a reason for caution.
